\section{Survey items derived from the inventory}
\label{sec:survey}

Table~\ref{tab:survey} sets out fourteen candidate items for a survey of
laboratory practice, derived from the inventory. Each is tied to the step it
interrogates and to the reason it is worth asking.

The selection principle is that a survey question earns its place when the
answer is actionable and currently unknown. Questions whose answers are
already predictable, or on which the evidence is strong and practice is
good, are omitted deliberately. That is why no item asks whether cartridge-based
testing is performed in a cabinet: the measured evidence makes that a low-yield
question.

Items are grouped by what they establish. Items 1--4 establish whether the
pre-analytical bench is treated as a respiratory-hazard area before anyone
knows what the specimen contains. Items 5--7 probe the failure modes that the
conditionality analysis identifies as carrying the risk. Items 8--10 address
the E5 gap at automated analysers, where the survey is partly substituting for
absent measurement. Items 11--13 test the controls on which the
strongest-evidence recommendations depend. Item 14 is a calibration question.

\begin{table*}[t]
\centering
\caption{Candidate survey items derived from the procedure inventory. Grade is
the evidence grade of the step the item interrogates.}
\label{tab:survey}
\small
\begin{tabular}{@{}p{0.025\textwidth} p{0.055\textwidth} p{0.035\textwidth} p{0.40\textwidth} p{0.40\textwidth}@{}}
\toprule
 & \textbf{Step} & \textbf{Gr.} & \textbf{Item} & \textbf{Why it is worth asking} \\
\midrule
1 & S01 & E2 & Where are respiratory specimens flagged for mycobacterial
investigation first opened, and is that location distinct from general
specimen reception? & Scores in the High band on the inhalational pathway, yet
is commonly performed on an open bench. \\

2 & S01 & E2 & Is there a written trigger for transferring a leaking or
damaged container into a cabinet unopened, and where is it displayed? & The
dropped-container accident was measured two orders of magnitude above
centrifuge rupture. \\

3 & S03 & E1 & Is closed-tube sampling configured and in routine use, as
distinct from being available on the instrument? & Removes the step rather
than controlling it; availability and use are commonly confused. \\

4 & S03 & E1 & Is decapping technique assessed for competency, and how
often? & The negative measurement was obtained under controlled removal; the
reassurance is technique-conditional. \\

5 & S02 & E1 & Are sealed safety buckets used, and are they opened inside a
cabinet for specimens under mycobacterial or fungal investigation? & The
negative centrifugation result was obtained with closed-cap tubes in a
swing-bucket rotor. \\

6 & S02 & E1 & Is there a documented inspection schedule for rotors, buckets,
gaskets and O-rings, with records retained? & The step is
failure-conditional; inspection is the control that addresses the failure. \\

7 & S02 & E1 & Is there a written breakage-in-the-centrifuge procedure
specifying who is notified and when the room is re-entered? & Frequently
absent; the routine run is not where the risk sits. \\

8 & S04 & E5 & Is clearing a probe jam or short sample treated as a distinct
task with its own written controls, or as part of running the analyser? &
Cover-open work is where release would occur; no measurement exists. \\

9 & S04 & E5 & Is liquid waste plumbed to a closed drain, or decanted from a
carboy by hand? & An uncharacterised release opportunity at the highest
throughput in the laboratory. \\

10 & S04 & E5 & Are manufacturer interlocks intact and unmodified? &
Defeated interlocks convert a closed pathway into an open one. \\

11 & S07 & E3 & Are all manipulations of sterile-site isolates performed in a
certified Class~II cabinet? & The control carried by the strongest evidence in
the inventory. \\

12 & S07 & E3 & When was each cabinet last field-certified, to what standard,
and are records retained? & The culture-step recommendation depends entirely
on the cabinet functioning. \\

13 & S06 & E3 & Is there a standing rule diverting slow-growing blood culture
isolates to a cabinet and to rule-out, irrespective of Gram stain appearance? &
In three of four documented events the stain was misleading; a stain-based
prompt would not have fired. \\

14 & --- & --- & Does the laboratory's own procedure-level risk assessment
record which of its controls rest on measurement and which on precaution? &
Calibration item: distinguishes a laboratory that has assessed its bench from
one that has documented it. \\
\bottomrule
\end{tabular}
\end{table*}

Two cautions apply to using these items.

They are candidate items, not a validated instrument. They have not been
cognitively tested, piloted, or assessed for response distribution, and the
phrasing will need adjustment to the surveying body's conventions.

They are also diagnostic rather than evaluative. An item such as the
certification question establishes whether a control is in place; it does not
establish whether the laboratory is safe, and a survey that scored laboratories
on these items without that caveat would be over-reading them.
